\section{Sharp examples and no-go models}
\label{sec:counterexamples}

The following examples separate exact determinant geometry from signed
cancellation.

\subsection{One anchor, both signs}

\begin{example}[Opposite signs with the same rows and anchor]
\label{ex:opposite-sign-rectangles}
Take
\[
 m=3,\qquad n=5,\qquad h_0=1,\qquad d_m=d_n=1.
\]
Then \(g=a=b=1\) and \(H=2\).  Choose \(j_R=2\).  The anchor targets are
\[
 A_m=7=7\cdot1,\qquad
 A_n=11=11\cdot1.
\]
Thus \(R_m=R_n=1\), the residuals agree, and \(t=1\).

First take \(j_E=4\).  Then
\[
 B_m=13=13\cdot1,\qquad
 B_n=21=7\cdot3.
\]
Hence
\[
 q=1,\qquad (u,v)=(1,3),
\]
and the rectangle sign is
\[
 \mu(1)\mu(1)\mu(1)\mu(3)=-1.
\]
The cross determinant is
\[
 7\cdot21-13\cdot11=4=H(j_E-j_R).
\]

Now take \(j_E=18\).  Then
\[
 B_m=55=11\cdot5,\qquad
 B_n=91=13\cdot7,
\]
where the displayed primes are the largest prime factors.  Thus
\[
 q=1,\qquad (u,v)=(5,7),
\]
and the sign is
\[
 \mu(5)\mu(7)=+1.
\]
Again all four targets are squarefree and
\[
 7\cdot91-55\cdot11=32=H(j_E-j_R).
\]
Therefore the four-corner identities and the divisor conditions admit
both signs even with the rows and same-residual anchor fixed.
\end{example}

\begin{remark}[Scope of the example]
\label{rem:small-example-scope}
The example realizes the exact squarefree largest-prime arithmetic
dictionary.  It is not asserted to satisfy every scale-dependent
missing-high, block, or endpoint mask of the full TPC carrier.  Its role
is sharp: no argument using only the identities proved in
\cref{sec:four-corner} can determine the sign.
\end{remark}

\subsection{Two points on one frozen ray}

\begin{example}[A nontrivial ray with constant sign]
\label{ex:frozen-ray-two-points}
Keep the same rows and anchor, and the erasure coordinates
\[
 (q,u,v)=(1,1,3).
\]
The point \(j_E=4\) above has
\[
 (Q_m,Q_n)=(13,7).
\]
\Cref{thm:erasure-ray} gives the candidate ray
\[
 j_E=4+3z,\qquad Q_m=13+9z,\qquad Q_n=7+5z.
\]
At \(z=2\), this gives
\[
 j_E=10,\qquad Q_m=31,\qquad Q_n=17,
\]
and indeed
\[
 3\cdot10+1=31\cdot1,\qquad
 5\cdot10+1=17\cdot3.
\]
Both \(z=0\) and \(z=2\) are squarefree largest-prime points, and both
have sign \(-1\).  The intermediate candidate \(z=1\) is deleted because
\(Q_m=22\) is not prime.  Thus primality thins the ray but does not change
its arithmetic sign.
\end{example}

\subsection{Kernel-level obstruction}

\begin{proposition}[No cancellation from the coordinate change alone]
\label{prop:kernel-level-no-go}
For every finite set
\[
 \Omega\subset
 \{(u,v)\in\N^2:(u,v)=1,\ (u,v)\ne(b,a)\}
\]
on which \(\mu(u)\mu(v)=\sigma\in\{-1,1\}\), there is a nonnegative kernel
\(K\) supported on \(\Omega\) for which
\begin{equation}
 \left|
 \sum_{(u,v)}\mu(u)\mu(v)K(u,v)
 \right|
 =
 \sum_{(u,v)}K(u,v).
 \label{eq:kernel-no-cancellation}
\end{equation}
\end{proposition}

\begin{proof}
Choose any nonzero nonnegative function \(K\) supported on \(\Omega\).
The signed sum is \(\sigma\sum K\), proving the identity.
\end{proof}

The proposition is an abstract sharpness model, not a claim about the
actual physical kernel.  It proves that divisor compression and primitive
coordinates alone cannot imply a signed saving.  A future theorem must use
additional regularity, distribution, or phase structure of the literal
kernel.

\subsection{Three common invalid inferences}

\begin{enumerate}[label=\textup{(\arabic*)},leftmargin=3em]
 \item From \(t,q\mid H\) one cannot infer that there are only
       divisor-many rectangles: the primitive pair \((u,v)\), the ray
       points, and the complete native labels remain.
 \item From \(t\mid H\) and \(q\mid H\) one cannot infer \(tq\mid H\).
       The exact cross relation only gives \(tq\mid H\Delta\).
 \item From two nonproportional prime forms one cannot infer signed
       cancellation.  Upper-bound sieves control their population, while
       the sign \(\mu(u)\mu(v)\) is constant on the whole ray.
\end{enumerate}
